% ===========================================================================
%  离散数学与结构（I） 2026-10-08（第 3--4 节）课堂笔记
%  群作用、类方程与 Sylow 定理
%  Compile:  latexmk -xelatex main      (or: xelatex main, twice)
% ===========================================================================
\documentclass[cjk]{loom}

\newcommand{\N}{\mathbb{N}}
\newcommand{\Z}{\mathbb{Z}}
\newcommand{\Q}{\mathbb{Q}}
\newcommand{\Stab}{\mathrm{Stab}}          % 稳定子
\newcommand{\Orb}{\mathrm{Orb}}            % 轨道
\newcommand{\Conj}{\mathrm{Conj}}          % 共轭类
\newcommand{\Syl}{\mathrm{Syl}}            % Sylow p-子群集合
\newcommand{\normal}{\trianglelefteq}      % 正规子群

\runningthread{离散数学与结构 · 2026-10-08}

\begin{document}

\loomcover
  {群作用、类方程与 Sylow 定理}
  {离散数学与结构（I）· 2026-10-08（第 3--4 节）课堂笔记}
  {}
  {2026 年 10 月}

\begin{strand}
本讲脉络：从群的分解纲领出发——阿贝尔群的结构已经清楚，非阿贝尔群则引出
单群与合成列；随后引入群作用、轨道与稳定子这一对基本工具；把轨道分解
运用到共轭作用上得到类方程，并立刻收获两个推论（$p$-群中心非平凡、
$p^2$ 阶群阿贝尔）；再由双重计数得到 Burnside 引理（轨道计数）；
接着讨论 $S_n$ 的共轭类与 $S_4$ 的类方程；最后证明 Sylow 三大定理，
并应用于 $pq$ 阶群的分类（及若干数例）。
\end{strand}

\clearpage
{\small\tableofcontents}

\clearpage
% ===========================================================================
\section{分解纲领：单群与合成列}\label{sec:structure}

\subsection{群的分解纲领：从阿贝尔群到非阿贝尔群}

\block{回顾（阿贝尔群）}
阿贝尔群可以分解为若干循环群的积；每个循环群又可分解为若干
\keyword{素数幂阶循环群}的积。所以任何有限阿贝尔群都可以分解成素数幂阶
循环群的积——对阿贝尔群而言，结构问题基本就清楚了。

\block{非阿贝尔群的纲领}
对非阿贝尔群 $G$，一般的目标是得到一列正规子群
\[
  \{e\} \normal G_1 \normal G_2 \normal G_3 \normal G,
\]
有了这列正规子群，就可以把问题分解成一串“因子”：
\[
  G_1,\quad G_2/G_1,\quad G_3/G_2,\quad G/G_3,
\]
即把群分解成若干\keyword{商群}分别研究。我们希望把它分解成“最简”的形式——
这就引出单群的概念。

\subsection{单群与合成列}

\begin{definition}[单群（Simple Group）]
群 $G$ 称为\keyword{单群}，如果它没有非平凡的正规子群，即
\[
  N \normal G \;\Longrightarrow\; N = G \ \text{或}\ N = \{e\}.
\]
\end{definition}

\begin{example}
素数阶循环群都是单群。
\end{example}

\begin{theorem}[合成列的存在性]
任何\keyword{有限群} $G$ 都存在一列如上的正规子群，使得每个因子
$G_{i+1}/G_i$ 都是单群。
\end{theorem}

\begin{proof}
假如某个因子不是单群，比如 $G/G_3$ 不是单群，那么它有一个非平凡正规子群
$\bar N$（既不等于 $G/G_3$，也不是平凡群）。由\keyword{商群的对应定理}，
$\bar N$ 在自然同态 $\pi: G\to G/G_3$ 下的\keyword{原像}
$\pi^{-1}(\bar N)$ 是 $G$ 的正规子群，且严格夹在 $G_3$ 与 $G$ 之间。
把这个中间群“切进”子群列中，得到更细的分解；有限群保证这一过程在有限步
后终止，归纳即得。
\end{proof}

% ===========================================================================
\section{群作用与轨道}\label{sec:actions}

\subsection{群作用}

\begin{definition}[群作用（Group Action）]
群 $G$ \keyword{作用}在集合 $X$ 上，即映射 $(g,x)\mapsto g\cdot x$，满足：
\begin{enumerate}[leftmargin=1.6em,itemsep=1pt]
  \item 单位元作用平凡：$e\cdot x = x$；
  \item 结合律：$g_1\cdot(g_2\cdot x) = (g_1 g_2)\cdot x$。
\end{enumerate}
\end{definition}

\begin{example}[共轭作用]
\leavevmode
\begin{enumerate}[leftmargin=1.6em,itemsep=2pt]
  \item \keyword{共轭作用}（conjugation）：群作用在自己身上，
        \[
          g\cdot x := gxg^{-1}.
        \]
  \item \keyword{在全体子群组成的集合上的共轭作用：}设
        $\mathcal{L}(G) = \{\,H \mid H\le G\,\}$（$G$ 的所有子群组成的集合），定义
        \[
          g\cdot H := gHg^{-1} = \{\,ghg^{-1} \mid h\in H\,\}.
        \]
        不难验证 $gHg^{-1}$ 仍是子群，且满足群作用的两条公理，故这确实是
        一个群作用。
\end{enumerate}
\end{example}

\subsection{轨道与稳定子}

\begin{definition}[轨道（Orbit）]
元素 $x$ 的\keyword{轨道}是所有通过群作用能到达的地方：
\[
  \Orb(x) = \{\,g\cdot x \mid g\in G\,\}.
\]
\end{definition}

\begin{definition}[稳定子（Stabilizer）]
\[
  \Stab(x) = \{\,g\in G \mid g\cdot x = x\,\},
\]
即所有使 $x$ 不动的群元素；不难验证 $\Stab(x)\le G$。
\end{definition}

\block{两条基本性质}
\begin{itemize}[leftmargin=1.3em,itemsep=2pt]
  \item 任何两个轨道要么相同、要么不交——轨道给出 $X$ 的一个划分，
        故 $|X|$ 等于各轨道大小之和；
  \item 下面的定理把轨道大小与稳定子大小联系起来。
\end{itemize}

\begin{theorem}[轨道与稳定子定理]
\[
  |\Orb(x)| = \frac{|G|}{|\Stab(x)|}.
\]
\end{theorem}

\begin{strand}
直观解释：想象 $x$ 所在的轨道是一堆点；$\Stab(x)$ 是让 $x$ 不动的群元素。
轨道中每个点的“到达方式”恰好都有 $|\Stab(x)|$ 种（到达同一点的群元素
相差一个稳定子中的元素），所以 $|\Orb(x)|\times|\Stab(x)| = |G|$。
\end{strand}

% ===========================================================================
\section{类方程及其应用}\label{sec:class-equation}

\subsection{类方程（Class Equation）}

\begin{definition}[中心（Centre）]
群 $G$ 的\keyword{中心}为
\[
  C_G = \{\,x\in G \mid \forall g\in G,\ gx = xg\,\}.
\]
\end{definition}

把一般的轨道分解等式
\[
  |X| = \sum_{\text{每个轨道取代表 } x_i} \frac{|G|}{|\Stab(x_i)|}
\]
应用到\keyword{共轭作用}（此时 $X = G$）。共轭作用下 $x$ 的轨道
\[
  \Conj(x) = \{\,gxg^{-1} \mid g\in G\,\}
\]
称为 $x$ 的\keyword{共轭类}（Conjugacy Class）。把轨道按大小分两类：

\begin{itemize}[leftmargin=1.3em,itemsep=2pt]
  \item \keyword{大小为 1 的轨道：}$|\Orb(x)| = 1$ 意味着对任何 $g\in G$
        都有 $gxg^{-1} = x$，即 $x$ 与所有元素交换——这些元素恰好组成中心
        $C_G$，共贡献 $|C_G|$；
  \item \keyword{大小大于 1 的轨道：}共轭作用下
        \[
          \Stab(x) = \{\,g\in G \mid gxg^{-1} = x\,\} = C_G(x),
        \]
        即 $x$ 的\keyword{中心化子}（Centralizer，所有与 $x$ 交换的元素
        组成的子群）。
\end{itemize}

\begin{theorem}[类方程]
\[
  |G| = |C_G| + \sum_{i}\, [\,G : C_G(x_i)\,],
\]
其中求和跑遍每个大小大于 1 的共轭类的一个代表元 $x_i$。
\end{theorem}

\subsection{类方程的应用}

\begin{definition}[$p$-群]
群 $G$ 称为 \keyword{$p$-群}，如果 $|G|$ 是素数 $p$ 的幂。
\end{definition}

\subsubsection{$p$-群的中心非平凡}

\begin{theorem}[$p$-群的中心非平凡]
任何非平凡 $p$-群 $G$ 的中心都满足 $|C_G| > 1$。
\end{theorem}

\begin{proof}
看类方程 $|G| = |C_G| + \sum_i [G:C_G(x_i)]$：
\begin{itemize}[leftmargin=1.3em,itemsep=1pt]
  \item 左边 $|G|$ 是 $p$ 的幂，故被 $p$ 整除；
  \item 每一项 $[G:C_G(x_i)]$ 是 $|G|$ 的因数且大于 1，而 $|G|$ 是 $p$
        的幂，所以每一项都被 $p$ 整除；
  \item 因此 $|C_G|$ 也被 $p$ 整除，即中心里至少有 $p$ 个元素。
\end{itemize}
\end{proof}

\subsubsection{$p^2$ 阶群是阿贝尔群}

\begin{theorem}[$p^2$ 阶群是阿贝尔群]
若 $|G| = p^2$（$p$ 为素数），则 $G$ 是阿贝尔群。
\end{theorem}

\begin{proof}
由上一定理，$p\mid |C_G|$；又 $C_G\le G$，由 Lagrange 定理
$|C_G|\in\{p,\,p^2\}$。若 $|C_G|=p^2$ 则 $G=C_G$，证完。只需排除
$|C_G|=p$：

假设 $|C_G| = p$，取 $x\notin C_G$，考虑中心化子 $C_G(x)$：
\begin{itemize}[leftmargin=1.3em,itemsep=1pt]
  \item $C_G\subseteq C_G(x)$（中心元素与一切元素交换，当然与 $x$ 交换）；
  \item $x\in C_G(x)$ 而 $x\notin C_G$，故 $|C_G(x)| \ge p+1$。
\end{itemize}
但 $C_G(x)\le G$ 的大小只能是 $1, p, p^2$；既含 $p+1$ 个元素，只能
$|C_G(x)| = p^2$，即 $x$ 与所有元素交换，$x\in C_G$——矛盾。故
$|C_G|=p^2$，$G$ 阿贝尔。
\end{proof}

% ===========================================================================
\section{Burnside 引理与共轭类}\label{sec:burnside}

\subsection{Burnside 引理（轨道计数）}

\block{问题}
数集合 $X$ 在群 $G$ 作用下有多少条轨道。

\block{记号}
对 $g\in G$，记 $g$ 的\keyword{不动点集}为
\[
  S_g = \{\,x\in X \mid g\cdot x = x\,\}.
\]

\begin{theorem}[Burnside 引理]
\[
  \text{轨道个数} = \frac{1}{|G|}\sum_{g\in G} |S_g|.
\]
\end{theorem}

\begin{proof}
（双重计数）考虑集合 $T = \{\,(g,x)\in G\times X \mid g\cdot x = x\,\}$。
\begin{itemize}[leftmargin=1.3em,itemsep=2pt]
  \item \keyword{按 $g$ 拆：}固定 $g$，这样的 $x$ 恰有 $|S_g|$ 个，故
        \[
          |T| = \sum_{g\in G}|S_g|;
        \]
  \item \keyword{按 $x$ 拆：}固定 $x$，这样的 $g$ 恰有 $|\Stab(x)|$ 个，故
        \[
          |T| = \sum_{x\in X}|\Stab(x)|.
        \]
        同一条轨道上的元素有共轭的稳定子：若 $y = g\cdot x$，则
        $\Stab(y) = g\,\Stab(x)\,g^{-1}$，大小相同。于是每条轨道
        $\mathcal{O}$ 贡献
        \[
          \sum_{x\in\mathcal{O}}|\Stab(x)|
          = |\mathcal{O}|\cdot|\Stab(x)|
          = |G|,
        \]
        其中最末一步用轨道与稳定子定理。把各条轨道的贡献相加，得
        $|T| = (\text{轨道个数})\cdot|G|$。
\end{itemize}
合起来：
\[
  \sum_{g\in G}|S_g| \;=\; |T| \;=\; (\text{轨道个数})\cdot|G|,
\]
两边同除以 $|G|$ 即得 Burnside 引理。
\end{proof}

\begin{remark}[用法]
想数轨道个数，就对每个群元素数它的不动点，加起来再除以 $|G|$。
\end{remark}

\subsection{例子：对称群 $S_n$ 的共轭类}

考虑 $S_n$ 对自身的共轭作用；共轭作用的轨道称为\keyword{共轭类}
（conjugacy class），即 $\Conj(x) = \{\,gxg^{-1} \mid g\in S_n\,\}$。

\block{共轭在置换上的效果}
以 3-轮换为例，
\[
  g\,(1\,2\,3)\,g^{-1} = \big(g(1)\ g(2)\ g(3)\big).
\]
推导：$g(1)$ 先被 $g^{-1}$ 映到 $1$，再被 $(1\,2\,3)$ 映到 $2$，再被 $g$
映到 $g(2)$；其余同理。

\block{结论}
$S_n$ 中两个置换共轭，\keyword{当且仅当}它们拆成不相交轮换后轮换的长度
对应相等（\keyword{轮换型}相同）。例如一个 3-轮换无论如何共轭也不会变成
2-轮换。

\block{枚举 $S_4$ 的共轭类（按轮换型）}
\begin{center}
\begin{tabularx}{0.62\linewidth}{Lr}
  \toprule
  轮换型 & 个数 \\
  \midrule
  恒等 $e$                          & $1$ \\
  对换 $(\cdot\,\cdot)$             & $\binom{4}{2} = 6$ \\
  3-轮换                            & $\binom{4}{3}\times 2 = 8$ \\
  双对换 $(\cdot\,\cdot)(\cdot\,\cdot)$ & $3$ \\
  4-轮换                            & $6$ \\
  \bottomrule
\end{tabularx}
\end{center}
合计 $1+6+8+3+6 = 24 = |S_4|$。

\block{$S_4$ 的类方程}
$24 = 1 + 6 + 8 + 3 + 6$。其中“中心部分”只有恒等元自己：
$C_{S_4} = \{e\}$。

% ===========================================================================
\section{Sylow 定理}\label{sec:sylow}

\subsection{Sylow 理论：定义}

设 $G$ 为有限群，$|G| = p^k m$，$p$ 为素数且 $p\nmid m$（$p^k$ 是 $|G|$
中 $p$ 的最高幂次）。

\begin{definition}[$p$-子群]
$H\le G$ 称为 \keyword{$p$-子群}，如果它是子群且 $|H|$ 是 $p$ 的幂。
\end{definition}

\begin{definition}[Sylow $p$-子群]
$H\le G$ 称为 \keyword{Sylow $p$-子群}，如果它是 $p$-子群且
$|H| = p^k$——$p$ 的幂次取到最大。
\end{definition}

\begin{example}[$G = S_4$]
$|S_4| = 24 = 2^3\cdot 3$。
\begin{itemize}[leftmargin=1.3em,itemsep=1pt]
  \item $p=3$：Sylow $3$-子群大小为 $3$，例如 $\langle(1\,2\,3)\rangle$；
  \item $p=2$：Sylow $2$-子群大小应为 $8$；其存在性由下面的 Sylow 第一
        定理保证，具体构造见本节末尾的例子。
\end{itemize}
\end{example}

\subsection{Sylow 第一定理（存在性）}

\begin{theorem}[Sylow 第一定理]
对任何有限群 $G$ 与素数 $p$，Sylow $p$-子群存在。
\end{theorem}

\begin{proof}
（对 $|G|$ 归纳）设 $|G| = p^k m$，$p\nmid m$，$k\ge 1$。归纳假设：
命题对所有更小的群成立。考虑类方程
\[
  |G| = |C_G| + \sum_i [G : C_G(x_i)],
\]
分两种情形。

\keyword{情形 1：$p\nmid |C_G|$。}左边被 $p$ 整除而 $|C_G|$ 不被 $p$
整除，故右端至少有一项不被 $p$ 整除：存在 $x_i$ 使 $p\nmid [G:C_G(x_i)]$，
从而 $p^k \mid |C_G(x_i)|$。又 $C_G(x_i)$ 比 $G$ 小（因为
$[G:C_G(x_i)]>1$），由归纳假设其中有 $p^k$ 阶子群，即 $G$ 的 Sylow
$p$-子群。

\keyword{情形 2：$p\mid |C_G|$。}中心是阿贝尔群，由（阿贝尔群的）
Cauchy 定理，中心里有 $p$ 阶元素 $c$。考虑 $\langle c\rangle$：
\begin{itemize}[leftmargin=1.3em,itemsep=2pt]
  \item $\langle c\rangle\le C_G$，其中元素与 $G$ 的一切元素交换，
        故 $\langle c\rangle\normal G$；
  \item 商群 $G/\langle c\rangle$ 的阶为 $p^{k-1}m$，由归纳假设其中有
        $p^{k-1}$ 阶子群 $\bar H$；
  \item 取 $\bar H$ 在自然同态 $\pi: G\to G/\langle c\rangle$ 下的原像
        $H=\pi^{-1}(\bar H)$，则 $H\le G$ 且
        $|H| = |\bar H|\cdot|\langle c\rangle| = p^{k-1}\cdot p = p^k$，
        即为 Sylow $p$-子群。
\end{itemize}
\end{proof}

\begin{remark}[回到 $S_4$]
先看 $\{e,\ (1\,2),\ (3\,4),\ (1\,2)(3\,4)\}$——一个 4 阶 $2$-子群
（还不是 Sylow）；再把 $(1\,3)(2\,4)$ 加进来，膨胀成一个 8 阶子群——
这就是 $S_4$ 的 Sylow $2$-子群（同构于二面体群 $D_4$，阶为 $8$）。
\end{remark}

\subsection{Sylow 第二定理}

\begin{theorem}[Sylow 第二定理]
设 $P$ 是 $G$ 的任一 Sylow $p$-子群，$Q$ 是 $G$ 的任一 $p$-子群，
则存在 $g\in G$ 使得
\[
  Q \le gPg^{-1},
\]
即任何 $p$-子群都含于该 Sylow $p$-子群的某个共轭之中。
\end{theorem}

\begin{corollary}
\leavevmode
\begin{enumerate}[leftmargin=1.6em,itemsep=2pt]
  \item \keyword{任意两个 Sylow $p$-子群互相共轭。}取 $Q$ 也是 Sylow
        $p$-子群，则 $Q\le gPg^{-1}$，而两边大小相同，故 $Q = gPg^{-1}$。
  \item \keyword{任何 $p$-子群都含于某个 Sylow $p$-子群。}
  \item \keyword{Sylow $p$-子群唯一 $\iff$ 它是正规子群。}
        \begin{itemize}[leftmargin=1.3em,itemsep=1pt]
          \item $(\Leftarrow)$ 若 $P\normal G$：对任何 $g$，
                $gPg^{-1} = P$；而所有 Sylow $p$-子群都是 $P$ 的共轭，
                故只有自己一个。
          \item $(\Rightarrow)$ 若 $P$ 唯一：对任何 $g$，$gPg^{-1}$ 仍是
                Sylow $p$-子群（共轭不改变大小），由唯一性 $gPg^{-1}=P$，
                故 $P$ 正规。
        \end{itemize}
\end{enumerate}
\end{corollary}

\begin{proof}[Sylow 第二定理的证明]
设 $P$ 为 Sylow $p$-子群，$Q$ 为任意 $p$-子群，让 $Q$ \keyword{左乘作用}
在 $P$ 的全体左陪集组成的集合 $G/P$ 上：
\[
  q\cdot(gP) := (qg)P.
\]
注意这里的 $G/P$ \keyword{不是商群}，只是左陪集的集合（没有假设 $P$
正规）。不难验证这确实是一个群作用。
\begin{itemize}[leftmargin=1.3em,itemsep=2pt]
  \item 陪集个数 $[G:P] = m$ 不被 $p$ 整除（$P$ 已把 $|G|$ 中 $p$ 的
        幂次全用掉了）。而 $[G:P]$ 等于各轨道大小之和，所以
        \keyword{必存在一条轨道，其大小不被 $p$ 整除}；
  \item 另一方面，由轨道与稳定子定理，每条轨道大小整除 $|Q|$（$p$ 的幂）；
        不被 $p$ 整除的 $p$ 的幂只有 $1$——故存在\keyword{不动点} $gP$：
        对任意 $q\in Q$ 有 $qgP = gP$。
\end{itemize}
两边左乘 $g^{-1}$ 得 $g^{-1}qg\in P$，即
\[
  Q \subseteq gPg^{-1}.
\]
\end{proof}

\subsection{正规化子（Normalizer）}

\begin{definition}[正规化子（Normalizer）]
设 $P\le G$，$P$ 在 $G$ 中的\keyword{正规化子}为
\[
  N_G(P) = \{\,g\in G \mid gPg^{-1} = P\,\},
\]
即共轭作用下使 $P$ 不变的元素全体（也就是子群集合上共轭作用中 $P$ 的
稳定子）。
\end{definition}

\begin{itemize}[leftmargin=1.3em,itemsep=1pt]
  \item 易验证 $N_G(P)\le G$，且显然 $P\subseteq N_G(P)$；
  \item 顾名思义，$P\normal N_G(P)$——我们只挑选了使 $P$ 在共轭下不变
        的元素。
\end{itemize}

\begin{lemma}
设 $P$ 是 $G$ 的 Sylow $p$-子群，则 $P$ 是 $N_G(P)$ 的\keyword{唯一}
Sylow $p$-子群。
\end{lemma}

\begin{proof}
由 $P\le N_G(P)\le G$ 知 $|N_G(P)|$ 中 $p$ 的最高幂次仍为 $|P|$，故
$P$ 是 $N_G(P)$ 的 Sylow $p$-子群；又 $P\normal N_G(P)$，把上面推论
第 3 条（唯一 $\iff$ 正规）应用于群 $N_G(P)$，即得 $P$ 是 $N_G(P)$ 的
唯一 Sylow $p$-子群。
\end{proof}

\subsection{例子：对称群的 Sylow 子群的构造}

\begin{example}[$S_4$ 的 Sylow $2$-子群（阶 8）]
由如下置换生成：
\begin{itemize}[leftmargin=1.3em,itemsep=1pt]
  \item $\{1,2\}$ 与 $\{3,4\}$ 两块\keyword{整体交换}：$(1\,3)(2\,4)$；
  \item $\{1,2\}$ 内部交换：$(1\,2)$；
  \item $\{3,4\}$ 内部交换：$(3\,4)$。
\end{itemize}
\end{example}

\begin{example}[$S_{p^2}$ 的 Sylow $p$-子群]
类似想法可构造一般的 $S_n$ 的 Sylow $p$-子群。以 $S_{p^2}$ 为例：
\begin{itemize}[leftmargin=1.3em,itemsep=2pt]
  \item 把 $1,\dots,p^2$ 分成 $p$ 组：$\{1,\dots,p\}$、
        $\{p+1,\dots,2p\}$、…、$\{(p-1)p+1,\dots,p^2\}$；
  \item $p$ 个组之间可以\keyword{整体轮换}（1 号组→2 号组→…）；
  \item 每组内部各有一个 $p$-轮换。
\end{itemize}
这样得到一个 $p^{p+1}$ 阶子群，正是 $S_{p^2}$ 的 Sylow $p$-子群。
$n$ 更大时构造类似。
\end{example}

\begin{remark}
由 Sylow 第二定理，Sylow $p$-子群彼此共轭，所以其它方式构造出的
Sylow $p$-子群都与上述构造共轭。
\end{remark}

\subsection{Sylow 第三定理}

\begin{theorem}[Sylow 第三定理]
设 $|G| = p^n m$，$p\nmid m$，记
\[
  n_p = \#\{\,\text{$G$ 的 Sylow $p$-子群}\,\}.
\]
则
\[
  n_p \mid m,\qquad n_p \equiv 1 \pmod p.
\]
\end{theorem}

\begin{proof}[$n_p\mid m$ 的证明]
考虑 $G$ 在全体 Sylow $p$-子群组成的集合 $\Syl_p(G)$ 上的
\keyword{共轭作用}：
\begin{itemize}[leftmargin=1.3em,itemsep=2pt]
  \item 由 Sylow 第二定理，任意两个 Sylow $p$-子群共轭，故该作用
        \keyword{只有一个轨道}——这种作用称为\keyword{传递作用}
        （Transitive Action，亦称可迁）；
  \item 由轨道与稳定子定理，轨道大小 $n_p = [G : \Stab(P)]$，而这里的
        稳定子正是正规化子：
        \[
          \Stab(P) = \{\,g \mid gPg^{-1} = P\,\} = N_G(P);
        \]
  \item 故 $n_p = [G : N_G(P)]$。由于 $P\le N_G(P)$，
        $n_p = [G:N_G(P)]$ 整除 $[G:P] = m$。
\end{itemize}
\end{proof}

\begin{proof}[$n_p\equiv 1\pmod p$ 的证明]
取定一个 Sylow $p$-子群 $P$，让 $P$ \keyword{共轭作用}在 $\Syl_p(G)$
上：
\[
  g\cdot P_i := gP_i g^{-1}\qquad (g\in P).
\]
把 $n_p$ 按轨道分解，每个轨道取代表 $P_i$：
\[
  n_p = \sum |\Orb(P_i)| = \sum [\,P : \Stab(P_i)\,],
\]
其中 $\Stab(P_i) = \{\,g\in P \mid gP_ig^{-1} = P_i\,\} \le P$。
分两种情形：
\begin{itemize}[leftmargin=1.3em,itemsep=2pt]
  \item \keyword{$\Stab(P_i) \ne P$：}此时 $[P:\Stab(P_i)]$ 是大于 1
        的 $p$ 的幂，故 $p \mid |\Orb(P_i)|$；
  \item \keyword{$\Stab(P_i) = P$：}即 $P$ 中每个元素都使 $P_i$ 在共轭
        下不变，故 $P\le N_G(P_i)$。由上一小节的引理，$P_i$ 是
        $N_G(P_i)$ 的\keyword{唯一} Sylow $p$-子群；而
        $P\le N_G(P_i)$ 且 $|P| = |P_i|$，$P$ 也是 $N_G(P_i)$ 的
        Sylow $p$-子群，所以 $P = P_i$。这条轨道就是 $\{P\}$，
        大小为 $1$。
\end{itemize}
于是只有一条大小为 1 的轨道（$P$ 自己），其余轨道大小都被 $p$ 整除，故
\[
  n_p \equiv 1 \pmod p.
\]
\end{proof}

% ===========================================================================
\section{Sylow 定理的应用}\label{sec:applications}

\subsection{应用举例}

\begin{example}[35 阶群]
设 $|G| = 35 = 5\cdot 7$。数 Sylow 子群的个数：$n_5\mid 7$ 且
$n_5\equiv 1\pmod 5 \Rightarrow n_5 = 1$；$n_7\mid 5$ 且
$n_7\equiv 1\pmod 7 \Rightarrow n_7 = 1$。唯一即正规：$G$ 有 5 阶正规
子群 $H_5$ 与 7 阶正规子群 $H_7$，于是 $G$ 是它们的\keyword{内直积}
（Internal Direct Product；数一下大小即知等于 $G$）：
\[
  G \cong H_5\times H_7 \cong C_5\times C_7 \cong \Z_{35}.
\]
即 35 阶群（在同构意义下只有循环群）。
\end{example}

\begin{example}[$11^2\cdot 13^2$ 阶群]
设 $|G| = 11^2\cdot 13^2$。与上例几乎一样：$n_{11}\mid 13^2$ 且
$n_{11}\equiv 1\pmod{11} \Rightarrow n_{11} = 1$，Sylow $11$-子群唯一
（故正规）；同理 Sylow $13$-子群唯一（故正规）。又由前面的定理，$p^2$
阶群必为阿贝尔群，故两个 Sylow 子群都阿贝尔。于是
\[
  G \cong P_{11}\times P_{13},
\]
两个阿贝尔群的直积仍阿贝尔，故 $G$ 阿贝尔。
\end{example}

类似方法适用于 $7\times 11$、$11\times 13$ 等阶数。

\subsection{$pq$ 阶群的分类}

\block{问题}
设 $|G| = pq$，$p<q$ 为素数。$G$ 有哪些可能？

\block{第一步：$n_q = 1$}
$n_q\mid p$ 且 $n_q\equiv 1\pmod q$；由 $q>p$ 知只能 $n_q = 1$。故 $G$
有唯一（因而正规）的 $q$ 阶子群，它是循环群，记
$Q = \langle g\rangle\normal G$；再取一个 Sylow $p$-子群
$H = \langle h\rangle$。

\block{第二步：$n_p\in\{1,\,q\}$}
$n_p\mid q$ 且 $n_p\equiv 1\pmod p$，故 $n_p = 1$ 或 $q$。分两种情形：
\begin{itemize}[leftmargin=1.3em,itemsep=2pt]
  \item \keyword{情形一：$q\not\equiv 1\pmod p$。}则 $n_p\ne q$，只能
        $n_p = 1$，回到上面数例的情形：$G\cong C_{pq}$ 循环。
  \item \keyword{情形二：$q\equiv 1\pmod p$。}此时 Sylow $p$-子群不一定
        唯一；$G$ 要么阿贝尔，要么同构于一个非阿贝尔群。
\end{itemize}

\block{分析（共轭给出的限制）}
因为 $Q = \langle g\rangle$ 正规，$hgh^{-1}\in\langle g\rangle$，故存在
$s$ 使
\[
  hgh^{-1} = g^{s}.
\]
对 $s$ 的限制：
\begin{itemize}[leftmargin=1.3em,itemsep=2pt]
  \item $s\ne 0$（显然）；
  \item \keyword{若 $s = 1$：}$h$ 与 $g$ 交换。群中每个元素都可写成
        $h^x g^y$；二者交换则整个群交换，$G$ 阿贝尔
        （$\cong C_p\times C_q$）。
\end{itemize}
把 $s$ 看作乘法群 $(\Z/q\Z)^\times$ 中的元素（该群阶为 $q-1$），则
$\mathrm{ord}(s)\mid q-1$；另一方面由 $h^p = e$ 得
\[
  g = h^p g h^{-p} = g^{s^p}\ \Longrightarrow\ s^p\equiv 1\pmod q,
\]
故 $\mathrm{ord}(s)\mid p$。于是
\[
  \mathrm{ord}(s)\mid \gcd(p,\,q-1).
\]
\begin{itemize}[leftmargin=1.3em,itemsep=2pt]
  \item 若 $p\nmid q-1$：立刻推出 $s = 1$，回到阿贝尔情形；
  \item 若 $p\mid q-1$：$s$ 可以取 $(\Z/q\Z)^\times$ 中阶为 $p$ 的元素。
\end{itemize}

\block{结论}
$s$ 的不同取值并不给出不同的群：通过适当更换生成元，这些取值彼此等价，
给出的群互相同构。因此：
\begin{itemize}[leftmargin=1.3em,itemsep=2pt]
  \item 若 $q\not\equiv 1\pmod p$，$pq$ 阶群只有循环群 $C_{pq}$；
  \item 若 $q\equiv 1\pmod p$，$pq$ 阶群除循环群外，（在同构意义下）
        恰有一个非阿贝尔群。
\end{itemize}


\end{document}
